\documentclass[a4paper,fontset=none]{ctexart}

\usepackage{anyfontsize,graphicx}
\usepackage{amsmath,amssymb,mathrsfs,wasysym}
\usepackage{autobreak}
\allowdisplaybreaks
\usepackage[T1]{fontenc}
\usepackage{textcomp}
\usepackage{amsthm}
\usepackage[libertine,vvarbb]{newtxmath}
\usepackage[scr=rsfso,frak=euler,bb=ams]{mathalfa}
\usepackage{bm}
\usepackage{indentfirst}
\setlength{\parindent}{0em}
\usepackage{parskip}
\usepackage{geometry}
\geometry{top=2cm,bottom=2cm,left=2cm,right=2cm}
\usepackage{fontspec}
\setmainfont{Libertinus Serif}
\setsansfont{Libertinus Sans}
\setmonofont{JetBrains Mono}
\setCJKmainfont{SimSun}[BoldFont=Noto Sans CJK SC Medium]

\begin{document}

\title{\textbf{物理宇宙学基础作业}}
\author{1900011604\ \ 张植竣\ \ 北京大学}
\date{2021年4月7日}
\maketitle

\begin{itemize}
    \item[1.(a)]
    该宇宙模型中的临界密度为：
    \[
    \rho_\mathrm{crit} = \frac{3H_0^2}{8\pi G} = 9.2045\times10^{-27}\,\mathrm{kg/m^3}
    \]
    由题意，重子物质中$75\%$为$^1\mathrm{H}$，$25\%$为$^4\mathrm{He}$，则当前时刻的电子数密度为：
    \begin{align*}
    n_0
    &= n_\mathrm{H} + 2n_\mathrm{He} \\
    &= 75\%\times\frac{\Omega_\mathrm{b}\rho_\mathrm{crit}}{m_\mathrm{H}} + 25\%\times\frac{\Omega_\mathrm{b}\rho_\mathrm{crit}}{m_\mathrm{He}} \\
    &= 0.1927\,\mathrm{m^{-3}}
    \end{align*}

    \item[1.(b)]
    只考虑电子对光子的Thomson散射，散射截面为：
    \[
    \sigma = \frac{1}{6\pi\varepsilon_0^2}\left(\frac{e^2}{m_\mathrm{e}c^2}\right)^2 = 6.6525\times10^{-29}\,\mathrm{m^2}
    \]
    光深：
    \begin{align*}
    \tau
    &= n\sigma l \\
    &= n_0(1+z)^3 \cdot \sigma \cdot \frac{r}{1+z} \\
    &= n_0\sigma r(1+z)^2
    \end{align*}
    考虑到：
    \[
    r = \frac{c}{H_0}\int_{0}^{z}\frac{\mathrm{d}z}{\sqrt{E(z)}} = \frac{c}{H_0}\int_{0}^{z}\frac{\mathrm{d}z}{\sqrt{\Omega_\mathrm{m}(1+z)^3 + \Omega_\Lambda}}
    \]
    则得光深关于红移的关系：
    \[
    \tau(z) = \frac{n_0\sigma c}{H_0}\int_{0}^{z}\frac{(1+z)^2\mathrm{d}z}{\sqrt{\Omega_\mathrm{m}(1+z)^3 + \Omega_\Lambda}}
    \]
    由$\tau_\mathrm{CMB}=0.09$及$\Omega_\mathrm{m}+\Omega_\Lambda=1$得：
    \begin{align*}
        \tau_\mathrm{CMB}
        &= \frac{n_0\sigma c}{H_0}\int_{0}^{z_\mathrm{reion}}\frac{(1+z)^2\mathrm{d}z}{\sqrt{\Omega_\mathrm{m}(1+z)^3 + \Omega_\Lambda}} \\
        &= \frac{n_0\sigma c}{H_0}\int_{1}^{x_\mathrm{reion}}\frac{\mathrm{d}x}{\sqrt{\Omega_\mathrm{m}x + \Omega_\Lambda}},\quad x = (1+z)^3 \\
        &= \frac{2n_0\sigma c}{3H_0\Omega_\mathrm{m}}(\sqrt{\Omega_\mathrm{m}x_\mathrm{reion}^3 + \Omega_\Lambda} - 1) \\
        &= \frac{2n_0\sigma c}{3H_0\Omega_\mathrm{m}}(\sqrt{\Omega_\mathrm{m}(1+z_\mathrm{reion})^3 + \Omega_\Lambda} - 1) \\
        &= 0.09
    \end{align*}
    计算得$z_\mathrm{reion}=11.33$，即宇宙再电离发生于约$z=11$处。
    
    \item[1.(c)]
    由光深的定义，观测者接收到红移为$z$处的辐射为：
    \[
    I = I_0\mathrm{e}^{-\tau}
    \]
    则光被星际介质散射的几率为：
    \[
    P = \frac{I-I_0}{I_0} = 1-\mathrm{e}^{-\tau}
    \]
    由$\tau=n_0\sigma r(1+z)^2$得：
    \[
    \frac{\mathrm{d}P}{R_0\mathrm{d}\omega} = \frac{\mathrm{d}P}{\mathrm{d}r} = \frac{\mathrm{d}P}{\mathrm{d}\tau}\cdot \frac{\mathrm{d}\tau}{\mathrm{d}r} = = n_0\sigma(1+z)^2\mathrm{e}^{-\tau(z)}
    \]
    
    \begin{figure}[!ht]
        \centering
        \includegraphics[width=0.75\linewidth]{Homework_3_3.png}
        \caption{$\mathrm{d}P/(R_0\,\mathrm{d}\omega) - z$关系图}
    \end{figure}

    \begin{figure}[!ht]
        \centering
        \includegraphics[width=0.75\linewidth]{Homework_3_4.png}
        \caption{$\mathrm{d}P/(R_0\,\mathrm{d}\omega) - D(z)$关系图，注意此处$D(z)=r$即为共动距离}
    \end{figure}
\end{itemize}

\end{document}